% Generated draft: requires mathematical and source review.
\documentclass[11pt]{article}
\usepackage[margin=25mm]{geometry}
\usepackage{fontspec}
\setmainfont{Latin Modern Roman}
\usepackage{amsmath,amssymb,mathtools,unicode-math}
\setmathfont{Latin Modern Math}
\usepackage{enumitem,needspace,xurl,hyperref}
\setlength{\parindent}{0pt}
\setlength{\parskip}{5pt}
\setlength{\emergencystretch}{3em}
\newcommand{\N}{\mathbb{N}}
\newcommand{\Z}{\mathbb{Z}}
\newcommand{\Q}{\mathbb{Q}}
\newcommand{\R}{\mathbb{R}}
\newcommand{\C}{\mathbb{C}}
\newcommand{\F}{\mathbb{F}}
\newcommand{\A}{\mathbb{A}}
\newcommand{\PP}{\mathbb{P}}
\newcommand{\E}{\mathbb{E}}
\newcommand{\eps}{\varepsilon}
\newcommand{\dd}{\,\mathrm{d}}
\begin{document}
\section*{UnsolvedMath ID: 30006426}
\textbf{Source problem code:} catalogue entry
\subsection{Definitions and mathematical statement}
This record collects several separate questions from an informal open-problem session. The original extraction contains two questions absent from the cleaned statement; both are retained below. Some necessary definitions and hypotheses are missing, so the collection does not specify a single fully determined mathematical problem.

\begin{enumerate}
\item The original extraction asks whether the continuous part of an induced action, transcribed as \(\Z\curvearrowright A_U\), lifts to \(\ell^\infty(A)\). For a \(C^*\)-algebra \(A\), \(\ell^\infty(A)\) denotes the \(C^*\)-algebra of norm-bounded sequences in \(A\). The source record does not specify \(A\), the meaning of \(A_U\), the inducing action, the meaning of its continuous part, or the quotient map and compatibility conditions defining a lift. Consequently, a precise lifting assertion cannot be recovered from the supplied text.

\item The original extraction asks whether
\[
 p\otimes 1_R\sim_{\mathrm{MvN}}1_R\otimes p
 \quad\text{in }R\otimes R.
\]
For projections \(e,f\) in a common \(C^*\)-algebra \(B\), Murray--von Neumann equivalence \(e\sim_{\mathrm{MvN}}f\) means that there exists \(v\in B\) with \(v^*v=e\) and \(vv^*=f\). Here the record does not define \(R\), specify the location or hypotheses on \(p\), or identify the tensor-product completion. Even the intended ambient category is therefore unresolved. The displayed expression is retained as a question with missing hypotheses, rather than asserted for all algebras or projections.

\item Let \(\mathcal O_\infty\) denote the Cuntz algebra on countably infinitely many canonical isometries, and let \(\mathcal K\) denote the compact operators on a separable infinite-dimensional Hilbert space. Put \(B=\mathcal O_\infty\otimes\mathcal K\). The supplied question asks whether there exists a \(*\)-automorphism \(\alpha:B\to B\) that fixes the diagonal and satisfies
\[
 K_0(\alpha)(x)=-x\qquad\text{for every }x\in K_0(B),
\]
where \(K_0(\alpha)\) is the homomorphism induced by \(\alpha\) on the operator \(K\)-theory group \(K_0(B)\). The diagonal subalgebra \(D\subseteq B\) is not identified. Moreover, ``fixes the diagonal'' might mean either
\[
 \alpha(D)=D
 \qquad\text{or}\qquad
 \alpha(d)=d\quad\text{for every }d\in D;
\]
the record does not resolve this distinction. The cleaned statement additionally asks whether such an automorphism can have the Rokhlin property. The original extraction separately asks, ``Does \(\alpha\) have the Rokhlin property?'' It does not explicitly establish the antecedent of this \(\alpha\), and no convention for the Rokhlin property on this nonunital algebra is supplied. Thus the cleaned reading is a possible grouping of the questions, with both its antecedent and the required Rokhlin convention requiring review.

\item Does there exist a relatively free group \(G\) such that \(C_r^*(G)\) is selfless? For a discrete group \(G\), its reduced group \(C^*\)-algebra \(C_r^*(G)\) is the norm closure in \(\mathcal B(\ell^2(G))\) of the linear span of its left regular representation, defined by \(\lambda_g\delta_h=\delta_{gh}\). Relative freeness requires an ambient variety of groups: a group is free in that variety when maps from its designated free generating set to any group in the variety extend uniquely to homomorphisms. The record does not specify the intended variety, rank, or additional restrictions on \(G\), and it does not define ``selfless.'' These missing choices prevent a fully specified existence assertion.

\item Does there exist a simple \(C^*\)-algebra \(A\) that is not \(K_1\)-injective? Here simple means nonzero with no nonzero proper closed two-sided ideals. For a unital \(C^*\)-algebra, \(K_1\)-injectivity means injectivity of the canonical map
\[
 U(A)/U_0(A)\longrightarrow K_1(A),
\]
where \(U(A)\) is its unitary group and \(U_0(A)\) is the component of the identity in the norm topology. Under this unital interpretation, failure of injectivity means that some \(u\in U(A)\) has zero \(K_1\)-class but does not belong to \(U_0(A)\). The supplied question does not impose unitality or separability, and does not state the convention for \(K_1\)-injectivity in the nonunital case. The unital explanation is therefore a definition for that case, not an added restriction on the question.

\item Is every \(\varepsilon\)-\(C^*\)-algebra close to a \(C^*\)-algebra? The cleaned statement replaces ``\(\varepsilon\)-\(C^*\)-algebra'' by ``approximate \(C^*\)-algebra.'' Neither expression is defined in the record. The permitted values of \(\varepsilon\), the approximate axioms, the meaning of closeness, and the required quantitative dependence are all unspecified. In particular, no assertion quantified over \(\varepsilon\) or over an approximation tolerance can be supplied without additional source information.
\end{enumerate}
\subsection{Short English statement}
Permanent UnsolvedMath ID \(30006426\) collects questions about lifting actions, equivalence of tensor-product projections, diagonal-fixing automorphisms, relatively free groups, \(K_1\)-injectivity, and perturbations of approximate \(C^*\)-algebras. Several essential hypotheses and definitions are missing. The supplied open classification concerns this unresolved extraction and is not an independently verified status assessment of each question.
\subsection{Sources}
\begin{itemize}
\item Local UnsolvedMath registry: supplied complete record with permanent ID \(30006426\), source problem code \texttt{catalogue entry}, titled ``Selected Lifting and Perturbation Problems in Operator Algebras.'' Repository registry locations are \texttt{data/unsolvedmath/problems.json} and the supplementary \texttt{catalog/problems.json}; their contents were not independently inspected for this draft.
\item Supplied source citation: \emph{C*-Algebras (2026)}, Oberwolfach Reports, DOI \texttt{10.4171/owr/2025/35}, pp.\ 1887--1948. Supplied source URL: \url{https://doi.org/10.4171/owr/2025/35}. The record describes these reports as documentation of mathematical workshops and their open-problem discussions.
\item Supplied attribution: ``Open Problem Session Collected by Stuart White (joint work with all and sundry).'' No individual proposer is supplied for these questions. The registry gives proposed year \(2025\), while the supplied report title contains \(2026\); this discrepancy is unresolved.
\item Supplied literature source: ``Oberwolfach Report 35,'' \url{https://doi.org/10.4171/owr/2025/35}.
\item Literature assessment dated exactly \texttt{2026-08-22}: supplied dated curation, not independently verified. It labels the record open because it aggregates unrelated questions and omits definitions and hypotheses, leaving no single determinate claim. The background calls this a dated literature triage derived from supplied OWR curation and requests independent verification of time-sensitive status and cited claims.
\item The registry labels its statement \texttt{corrected\_verified} and reports that the official report was checked and notation repaired. These are supplied provenance claims; this draft does not claim to have visited or independently verified the report.
\end{itemize}
\end{document}
